\chapter{Background}
\label{chap:background}

\section{Numerical variability}

As neuroscience has transitioned into an increasingly computational field, it has encountered challenges related to numerical reproducibility similar to those observed in many other scientific domains~\cite{Kiar2020}. In neuroimaging, where analyses must be re-executed to reproduce findings and evaluate their validity, numerous factors can contribute to variability between results. Data noise, software and hardware updates, operating-system differences, mathematical libraries, and parallelization strategies can introduce small numerical perturbations that produce deviations in the results and potentially compromise the reliability of computational analyses~\cite{Glatard2015,Salari2021}. Moreover, operations commonly employed in neuroimaging, including registration, segmentation, interpolation, and brain-network modelling, can be unstable when subjected to minor perturbations in the data or analytical pipeline~\cite{Kiar2020}. Numerical instability is therefore a fundamental consideration in computational reproducibility, and evaluating the resulting numerical variability is important for establishing the reliability of neuroimaging results.

\subsection{Sources of numerical variability}

Numerical errors may occur when non-integer data are represented using the IEEE~754 floating-point standard, which is widely employed to represent real numbers and perform arithmetic operations in computer systems. A floating-point number consists of a sign, an exponent, and a fraction. In single precision (32-bit), one bit is allocated to the sign, eight bits to the exponent, and 23 bits to the fraction; in double precision (64-bit), the corresponding allocations are one, 11, and 52 bits. Although this representation supports a broad range of very small and very large values, many real numbers have non-terminating binary representations and must therefore be approximated within the available number of bits~\cite{Salari2021}.

Two common consequences of finite-precision arithmetic are rounding error and catastrophic cancellation. Rounding error occurs when the result of an operation is rounded to fit within the available precision, producing a small deviation from the exact mathematical value. Catastrophic cancellation occurs when two nearly equal quantities are subtracted, causing a substantial loss of significant digits. These errors may accumulate across successive operations and can affect the accuracy and stability of a computation, particularly in long or nonlinear pipelines. Consequently, theoretically equivalent implementations may produce different results because of differences in operation order, compiler optimization, hardware architecture, or execution environment~\cite{Glatard2015,Kiar2020}.

\subsection{Assessment of numerical variability}

Stochastic arithmetic, uncertainty or sensitivity analysis, and random-seed analysis are among the methods available for detecting numerical instability. This thesis employs stochastic arithmetic, an approach used to investigate the effects of finite-precision arithmetic by modelling and estimating losses of numerical accuracy. A widely used technique within stochastic arithmetic is Monte Carlo Arithmetic (MCA), which models floating-point error through controlled random perturbations~\cite{Parker1997MCA,Salari2021}. For a floating-point value $x$, MCA defines an inexact counterpart as

\begin{equation}
  \operatorname{inexact}(x,t)=x+2^{e_x-t}\xi,
  \label{eq:mca-general}
\end{equation}

where $e_x$ is the exponent of $x$, $t$ is the virtual precision, and $\xi$ is a uniformly distributed random variable on $(-1/2,1/2)$. MCA can be applied in three perturbation modes. Precision Bounding perturbs the inputs of an operation and enables the detection of catastrophic cancellation; Random Rounding perturbs the output and enables the investigation of rounding error; and Full MCA combines both modes~\cite{Salari2021}.

Under the MCA framework, floating-point operations in an algorithm or software pipeline are randomly perturbed, and the analysis is repeated to generate a distribution of equally plausible results. The variability of this distribution provides information about the numerical stability of the evaluated software or pipeline. Numerical uncertainty may be summarized through the dispersion of repeated results, the number of significant digits that remain consistent across executions, or outcome-specific measures that compare numerical variability with other relevant sources of variation. The resulting estimates depend on the perturbation mode, the virtual precision, and the extent of software instrumentation~\cite{Denis2016Verificarlo,Kiar2021}. The particular MCA configuration and measures used in this thesis are presented in Chapter~\ref{chap:phase-one}.

\section{Numerical variability in neuroimaging: literature review}

Early evidence of numerical variability in neuroimaging arose from direct comparisons of computational environments. Glatard et al.~\cite{Glatard2015} evaluated FSL, FreeSurfer, and CIVET across different software builds and GNU/Linux operating systems. Although many outputs were similar, substantial differences occurred in particular analyses: Dice coefficients for FSL tissue classifications reached values as low as 0.59, independent-component analyses of fMRI data differed across operating systems for approximately one-third of the evaluated participants, and regional cortical-thickness estimates varied with the software build or operating system. The fMRI differences were traced to variations in motion-correction outputs. This study established that computational infrastructure can affect scientifically relevant neuroimaging results and emphasized the need to evaluate the numerical stability of complete pipelines.

Subsequent work used controlled perturbations to characterize this variability rather than relying exclusively on comparisons between existing environments. Kiar et al.~\cite{Kiar2021} instrumented deterministic and probabilistic structural-connectome pipelines with MCA and evaluated the resulting connectomes and network features. Stability ranged from nearly complete agreement to only zero or one significant digit. Although some aggregate properties of network organization remained stable, individual edge weights and participant-level features could be unreliable, demonstrating that group-level robustness does not necessarily imply reliability at the individual level. Salari et al.~\cite{Salari2021} subsequently introduced fuzzy libmath to perturb calls to mathematical libraries and reproduce variability caused by operating-system updates in Human Connectome Project preprocessing pipelines. The simulated and observed operating-system variability were comparable, while between-participant differences were preserved. However, the processed images retained fewer than eight significant bits of the 24 available, providing further evidence that infrastructure-level perturbations can propagate through neuroimaging workflows.

Numerical variability has also been investigated as a source of potentially useful variation. Using MCA-perturbed structural connectomes, Kiar et al.~\cite{Kiar2021Augmentation} constructed augmented datasets for age classification. Resampling across plausible perturbed networks improved test performance and generalizability across the evaluated classifiers, preprocessing methods, and dimensionality-reduction techniques. The improvement did not require a large number of perturbations. These results showed that numerical variability is not only a threat to reproducibility; when it is explicitly characterized, it can also be used to reduce computational bias and improve the robustness of downstream models.

Another line of work has incorporated numerical variability into software validation. Chatelain et al.~\cite{chatelain2024numerical} defined acceptable bounds of result variation from Random Rounding samples and used these bounds to construct stability tests for fMRIPrep. The tests accepted variations produced by a reference software version while detecting subtle image-processing changes and corrupted templates. This approach extends conventional software testing by recognizing that exact bitwise equality is not always an appropriate criterion for scientific pipelines. It also provides a mechanism for evaluating whether results obtained after software or hardware changes remain within the expected numerical variability of the reference analysis.

Recent studies have extended numerical-stability analysis to deep-learning-based neuroimaging. Gonzalez-Pepe et al.~\cite{Gonzalez2023Inference} applied Random Rounding during inference with SynthMorph and FastSurfer and compared the results with the conventional FreeSurfer pipeline. Across 32 participants, CNN inference was more numerically stable: nonlinear-registration outputs retained approximately 19 significant bits compared with 13 for the reference pipeline, and whole-brain segmentations achieved mean Dice scores of 0.99 compared with 0.92. In contrast, their later investigation of FastSurfer training found substantial training-time variability, particularly in cortical regions, which exceeded that of FreeSurfer~\cite{Gonzalez2025CNN}. Random-seed variability exhibited a spatial structure similar to numerical variability and could be used to construct ensembles that improved brain-age regression. Together, these findings distinguish the comparatively stable inference stage from the more variable iterative training stage.

Numerical variability can also depend strongly on the algorithm and its configuration. Mirhakimi et al.~\cite{mirhakimi2025numerical} compared linear registration with SPM, FSL, and ANTs across similarity measures, templates, and cohorts of healthy participants and participants with PD. SPM was the most stable under its default configuration, whereas FSL and ANTs exhibited greater variability, and ANTs occasionally produced registration failures under perturbation. Stability did not differ significantly between the healthy and PD cohorts, suggesting that evaluations conducted in healthy populations may generalize to clinical samples for this task. The study also demonstrated that numerical-variability measures may support automated quality control, connecting computational stability with the detection of problematic registrations.

Most recently, Chatelain et al.~\cite{Yohan2025NPVR} quantified the practical impact of numerical variability on structural MRI measures in PD using the Numerical--Population Variability Ratio. For several cortical and subcortical outcomes, numerical variability approached one-third of between-participant variability and altered conclusions concerning group differences and clinical associations. Propagation of numerical variability to 13 previously published PD studies produced an estimated mean probability of a significance-classification change of approximately 5\% for cross-sectional analyses and 10\% for longitudinal analyses. This work moves beyond reporting numerical differences by relating them directly to population variability and inferential decisions.

Taken together, these studies demonstrate that numerical variability can arise across operating systems, mathematical libraries, conventional image-processing algorithms, connectome pipelines, and deep-learning models. Its magnitude is specific to the pipeline, outcome, participant, and analytical configuration, and stability at the group level does not guarantee stability of regional or participant-level measures. Nevertheless, most previous investigations have focused on structural or diffusion MRI, image registration, segmentation, or structural connectivity. The numerical variability of functional-connectivity graph measures remains comparatively underexplored, motivating the Phase I study presented in Chapter~\ref{chap:phase-one}.
